\section{Method}
\label{sec:method}

\subsection{Agent and environment}
\label{sec:agent}

The agent is a LangGraph state-graph ReAct loop
\citep{yao2023react, langgraph2024}: \texttt{START} $\rightarrow$ agent
$\rightarrow$ tools $\rightarrow$ agent $\rightarrow \cdots \rightarrow$
\texttt{END}. At each turn the agent receives a text rendering of the page's
accessibility tree and emits either a tool call or a final answer. The tool set
comprises \ToolCount{} WebArena actions (click, type, hover, scroll, go-back,
go-forward, tab-focus, new-tab, and a no-op wait) dispatched over Playwright.

Two counter conventions matter for interpreting the cost columns later, so we fix
them now. \texttt{step\_count} counts \emph{executor} ReAct turns only;
planner and replanner invocations are counted in \texttt{llm\_calls} but not in
\texttt{steps}. In the ReAct configuration studied here the two coincide, but
they diverge under plan-and-execute, which is one reason the architecture
comparison in \secref{sec:results} is reported as a pilot rather than a ranking.

\subsection{Fault model}
\label{sec:faults}

We inject five faults, chosen to span the observation, network, and action
layers while remaining non-adversarial:

\begin{itemize}
  \item \textbf{DOM omission} (\texttt{web\_dom\_missing}) --- elements are dropped
  from the accessibility tree, so the agent observes a page that does not match
  the one it is acting on (observation layer).
  \item \textbf{Popup block} (\texttt{web\_popup\_block}) --- a blocking overlay
  occludes actionable content (observation layer).
  \item \textbf{HTTP error} (\texttt{web\_http\_error}) --- navigations return a
  server error (network layer).
  \item \textbf{Timeout} (\texttt{web\_timeout}) --- requests stall to the client
  deadline (network layer).
  \item \textbf{Parameter error} (\texttt{agent\_param\_error}) --- a tool call is
  executed with a corrupted argument, so it fails at the action layer and the
  agent must recover.
\end{itemize}

Each fault is injected at a fixed step index through a seeded
\texttt{FaultInjector}. The seed and the injection step are recorded per trial, so
a fault condition is reproducible exactly. Injection intensity is held at
\FaultIntensity{} throughout. The implementation contains ten fault
definitions in total; the five above are the ones with committed official
evaluations, and we restrict every empirical claim in this paper to those five.
The remaining five (targeting GitLab-specific state) are implemented but not part
of the main result, and we do not claim anything about them.

\subsection{Paired design and attribution quantities}
\label{sec:attribution}

For a fault $F$ and a (task, seed) cell $u$, let
$Y_u^{(0)}, Y_u^{(1)} \in \{0,1\}$ be official success in the control and fault
arms. The design runs both arms of every cell, so the unit of analysis is the
pair, and the estimand is the paired risk difference

\begin{equation}
  D(F) \;=\; \mathbb{P}\bigl(Y^{(1)} = 1\bigr) - \mathbb{P}\bigl(Y^{(0)} = 1\bigr),
\end{equation}

estimated by the mean of the within-pair differences
$d_u = Y_u^{(1)} - Y_u^{(0)}$. Because $d_u \in \{-1,0,+1\}$, the sample mean is
exactly the difference of the two marginal rates on the paired subset, and its
sign is determined by the two discordant counts $b$ (control passes, fault fails)
and $c$ (fault passes, control fails). We therefore report $b$, $c$, and an exact
McNemar test on them.

We deliberately do \emph{not} define a ``responsibility share'' between model and
architecture. Doing so requires a design in which both factors vary ---
$\mathrm{Model} \times \mathrm{Architecture} \times \mathrm{Fault}$ --- and ours
varies only the fault. A variance decomposition applied to these data would
attribute to ``model'' a quantity that is constant by construction.
\secref{sec:discussion} specifies the design that would identify it.

\subsection{Statistical procedure}
\label{sec:stats}

All tests are implemented in the Python standard library, with no SciPy
dependency, and are unit-tested against hand-computed values (see the
Reproducibility statement).

This section states the procedure in outline; \appref{app:stats} gives the full
derivations, the seed-sensitivity envelope, and the power table.

The primary test is the two-sided \emph{exact} McNemar test on the discordant
counts, reported with its mid-$p$ variant and a Haldane--Anscombe corrected
matched odds ratio, and with Holm--Bonferroni multiplicity control across the five
faults. Intervals for $D(F)$ are percentile bootstraps that resample \emph{pair
units} --- the correct unit for a paired design --- under a fixed seed, with a
task-clustered sensitivity interval and a multi-seed envelope as checks. Process
measures use the two-sided Wilcoxon signed-rank test with the Hodges--Lehmann
shift estimator, since the deltas are heavy-tailed and tied. Power is obtained by
solving the paired-proportion equation for $\delta$ at the \emph{observed}
discordance $\pi_d = \PiDNative{}$; because the risk difference is bounded by
$\pi_d$ rather than by $1$, the search must be bracketed on $(0, \pi_d]$, and
targets that are unreachable at any effect size are reported as \emph{undefined}
rather than as a number.
